% Generated draft: requires mathematical and source review.
\documentclass[11pt]{article}
\usepackage[margin=25mm]{geometry}
\usepackage{fontspec}
\setmainfont{Latin Modern Roman}
\usepackage{amsmath,amssymb,mathtools,unicode-math}
\setmathfont{Latin Modern Math}
\usepackage{enumitem,needspace,xurl,hyperref}
\setlength{\parindent}{0pt}
\setlength{\parskip}{5pt}
\setlength{\emergencystretch}{3em}
\newcommand{\N}{\mathbb{N}}
\newcommand{\Z}{\mathbb{Z}}
\newcommand{\Q}{\mathbb{Q}}
\newcommand{\R}{\mathbb{R}}
\newcommand{\C}{\mathbb{C}}
\newcommand{\F}{\mathbb{F}}
\newcommand{\A}{\mathbb{A}}
\newcommand{\PP}{\mathbb{P}}
\newcommand{\E}{\mathbb{E}}
\newcommand{\eps}{\varepsilon}
\newcommand{\dd}{\,\mathrm{d}}
\begin{document}
\section*{UnsolvedMath ID: 20003331}
\textbf{Source problem code:} catalogue entry
\subsection{Definitions and mathematical statement}
This record crosses a section boundary. Although its extracted first sentence reads as if Conjecture 39 were proved for some higher-rank symmetric spaces, the supplied aid says the preceding record supplies the words ``Results toward''. The qualified assertion is therefore that L. Silberman and A. Venkatesh proved results \emph{toward} Conjecture 39, cited as [178]. The aid identifies [178] as a nondegenerate Cartan-invariant microlocal-lift theorem, not full arithmetic quantum unique ergodicity. It distinguishes a later full theorem requiring compact prime-degree division-algebra spaces, Hecke eigenfunctions, and nondegenerate spectral hypotheses. Neither theorem's complete statement is supplied, so these qualifications must not be dropped.

Section 6.2 then begins a different discussion: M. Berry [19] models eigenfunctions of typical chaotic systems as superpositions of plane waves with random amplitudes, phases and directions, motivating Gaussian and independence predictions for eigenfunctions $\phi_i$ as $i\to\infty$, and a following conjecture for generic negatively curved compact Riemannian manifolds. That following conjecture is absent. No probability sampling, normalization, topology of ``generic'', or precise sense of independence is specified, so the prose is not a complete Gaussian-limit statement.

The supplied companion joining criterion, separate from either missing conjecture, has the following precise form. Let $X,Y$ be compact metric spaces, $\lambda_i$ probability measures on $X\times Y$, and suppose their two marginals converge weakly to probabilities $\mu,\nu$. For continuous real functions $f,g$ set
\[
 f_0=f-\int_X f\,d\mu,\qquad g_0=g-\int_Y g\,d\nu.
\]
Then
\[
 \lambda_i\longrightarrow\mu\otimes\nu
 \quad\Longleftrightarrow\quad
 \int_{X\times Y} f_0(x)g_0(y)\,d\lambda_i(x,y)\longrightarrow0
 \quad\text{for every }f\in C(X),\ g\in C(Y).
\]
Marginal convergence alone is insufficient: normalized Lebesgue measure on the diagonal of $[0,1]^2$ has uniform marginals but correlation $1/12$ for $f(x)=x-1/2$, $g(y)=y-1/2$.

For finite centered families $f_1,\ldots,f_r$ and $g_1,\ldots,g_s$, define the real matrix $C_{jk}=\int f_j(x)g_k(y)\,d\lambda(x,y)$. Its detector strength is
\[
 \sup_{\|u\|_2=\|v\|_2=1}
 \left|\int\Big(\sum_j u_jf_j\Big)\Big(\sum_kv_kg_k\Big)\,d\lambda\right|
 =\|C\|_{\mathrm{op}}.
\]
If $\lambda=c(\mu\otimes\nu)+(1-c)\rho$ with $0\leq c\leq1$ and $\rho$ a probability, a centered correlation defect $\delta=|\int f_0g_0\,d\lambda|>0$ implies
\[
 c\leq1-\frac{\delta}{\|f_0\|_\infty\|g_0\|_\infty}.
\]
When $\mu,\nu$ are Haar probabilities this bounds the product-Haar component. No assertion is made that such defects occur for the actual arithmetic or random-wave families in the fragment.
\subsection{Short English statement}
The entry joins qualified progress toward higher-rank arithmetic QUE to a separate introduction to Berry's random-wave predictions. Neither surrounding conjecture is fully stated. Its supplied companion criterion explains that convergent marginals imply joint independence only when all centered product correlations vanish.
\subsection{Sources}
\begin{itemize}
\item Local UnsolvedMath registry, permanent ID 20003331, exact source code catalogue entry; \texttt{aim-other-notes.json}, notes[137], end of the Conjecture 39 discussion and beginning of Section 6.2.
\item American Institute of Mathematics, \emph{Emerging applications of measure rigidity}, \url{https://aimath.org/WWN/measrigid/measrigid.pdf}. Numerical references [178] (L. Silberman and A. Venkatesh) and [19] (M. Berry) are retained; full bibliography entries are absent.
\item Supplied research attempt 1: context\_only, problem status not\_a\_problem, rigor label proved, novelty confidence low. It supplies the missing ``Results toward'' qualification, the microlocal-lift/later-theorem distinction, and the joining criterion and detectors. Machine-generated research aid, not peer reviewed; independently verify before citation. No complete citation to the later full theorem is supplied.
\end{itemize}
\end{document}
